% =====================================================================
%  starter-wtm.tex — copy this to start a "What to memorise" sheet.
%
%  The same style in a different shape: every quotation macro works here.
%  This file only shows what is DIFFERENT about a WTM sheet.
%
%  BUILD — from the kit's main folder:   bash tools/build.sh starter-wtm
%
%  THE ONE RULE THAT IS NOT TYPOGRAPHIC
%    Nothing here inserts content you did not type — no honorific, no
%    grading, no translation. If you want it on the page, type it.
% =====================================================================
\documentclass[12pt]{article}
\newcommand{\ibfontpath}{fonts/}
\usepackage[wtm]{islamicbooklet}

% Optional — leave out for a plain sheet:
% \bookletdomain{www.example.org}   % centred in the header
% \bookletlogo{my-logo.png}         % footer, bottom right

\begin{document}

% \wtmtitle{lesson code}{subject} — "What to memorise:" is fixed.
\wtmtitle{L1}{The Qualities of A Righteous Wife}

% Everything else sits inside wtmbody, which holds the sheet's 0.5in
% indent — Arabic paragraphs included.
\begin{wtmbody}

\wtmbullet{This document contains two (2) points}

% \wtmpoint{name} -> "1. name:"   The colon is added for you; an optional
% argument becomes a smaller parenthetical: \wtmpoint[optional]{...}
\wtmpoint{The Āyah}

% Under [wtm] the labels come out as "[meaning:]" and "Sūrah"; a booklet
% spells them "[meaning]:" and "Surah".
\ayah{فَالصَّالِحَاتُ قَانِتَاتٌ حَافِظَاتٌ لِلْغَيْبِ بِمَا حَفِظَ اللَّهُ}%
     {Therefore the righteous women are devoutly obedient \g{to Allâh
      and to their husbands}, and guard in the husband's absence what
      Allâh orders them to guard \g{e.g. their chastity, their
      husband's property}.}%
     {An-Nisāʾ}{4:34}

\wtmpoint[optional]{The Ḥadīth}

% On a WTM sheet the Arabic ḥadīth is its own paragraph — isnād,
% quotation and takhrīj in one RTL run, only the quotation green:
%   \arpara{...}   an RTL paragraph, black by default
%   \inhadith{...} turns a stretch of it green  (\inathar for an athar,
%                  \inquran for an āyah quoted inside a narration)
\arpara{عَنْ أَبِي هُرَيْرَةَ \ra{} قَالَ قَالَ رَسُولُ اللَّهِ \saw{}:
  \inhadith{«إِذَا صَلَّتِ الْمَرْأَةُ خَمْسَهَا وَصَامَتْ شَهْرَهَا وَحَصَّنَتْ فَرْجَهَا
  وَأَطَاعَتْ بَعْلَهَا دَخَلَتْ مِنْ أَيِّ أَبْوَابِ الْجَنَّةِ شَاءَتْ»}.
  رواه ابن حبان في صحيحه (٤١٦٣)، وحسنه الألباني.}

On the authority of Abū Hurayrah \hon{may Allah be pleased with him} who
said, The Messenger of Allah \saw{} said: \inhadith{“If a woman prays her
five prayers, fasts her month of Ramaḍān, guards her chastity, and obeys
her husband, she will enter Paradise from any gate she wishes.”} Narrated
by Ibn Ḥibbān in his \tr{Ṣaḥīḥ} (4163), declared \tr{ḥasan} by
Al-Albānī.

\end{wtmbody}

% ---------------------------------------------------------------------
%  NOTES
%  More points: add more \wtmpoint{...}. The header and footer repeat on
%  a second page by themselves. A WTM section inside a booklet: drop the
%  [wtm] option and call \wtmstyle where it begins (\bookstyle to end).
%  Transliteration: type ā ī ū ḥ ṣ ṭ ḍ ẓ ʿ ʾ directly; \ayn{} and
%  \hamzah{} exist for the two half rings. Never an apostrophe for either.
% ---------------------------------------------------------------------

\end{document}
